\section{Designs for a Routing Check}
\label{sec:routing-check-designs}

To attribute a wrong decision to the coordinator, and to prevent the wrong routing from taking effect, a check needs to be available before the coordinator invokes the selected sub-agent.

One option is a deterministic check, which applies a heuristic or a policy, encoded as a rule, over the request. Section~\ref{subsec:building-approaches} makes the same point about a hand-authored specification in general, in the context of classical agents. Writing it correctly demands domain expertise, and it is incomplete by construction, since it covers only the cases it was written to anticipate. A rule matched against a request is this kind of specification, so its coverage is bounded the same way, and it does not generalize across a sub-agent catalog it was not written for.

A language model can also serve as this check, reused as a judge that grades the coordinator's routing decision. As Section~\ref{sec:llm-as-judge} shows, a judge model carries systematic biases of its own. It also adds latency and cost, since it requires a separate model call for every decision it checks. Moreover, because the judge is itself an LLM performing a similar reasoning task, it is subject to the same class of failure as the coordinator it is meant to check.

A classifier trained to predict the correct sub-agent label from a fixed representation of the request's text can be a further source for evaluation. The representation is computed without the routing labels, and only the classifier that reads it is trained on them. This work refers to a classifier of this kind as a \emph{text-only classifier}. It can be trained on lexical features, counting how many times each word occurs in a query (a \emph{bag-of-words} representation), optionally weighted by how rare the word is across the training set using \emph{TF-IDF}. An alternative is a classifier trained on \emph{sentence embeddings}: dense vectors produced by an embedding model that capture the meaning of a sentence. In this case, two queries phrased differently but with the same meaning can map to nearby vectors. A text-only classifier of any of these three kinds can be more efficient than an LLM, and can also exceed a zero-shot LLM's accuracy on the same routing task \cite{balaMultiAgentRoutingSetValued2026}. However, this is not guaranteed: depending on the routing task, an LLM may perform better than a text-only classifier. Where a text-only classifier is more accurate than the coordinator's own LLM, its disagreement with the coordinator's decision could flag a likely misrouting, or the classifier could replace the coordinator's classification entirely, rather than only check it.

Wu, Zhu et al. showed that a probe, defined in Section~\ref{subsec:probing-method}, can capture a decision of this kind directly from a model's hidden states, with high accuracy: the tool a model is about to call is linearly encoded in its hidden states, and readable before the call is generated (Section~\ref{subsec:readability}) \cite{wuToolCallingLinearly2026}. A routing decision selects one sub-agent from a declared set, the same kind of selection that Wu, Zhu et al. studied for tool calls, so the information a check needs may already be present in the coordinator's hidden states before it generates a response. A probe trained on these hidden states could therefore be a stronger source than a text-only classifier for this decision, classifying the correct sub-agent directly from them and serving as an evaluation flag before the coordinator's decision takes place. However, this probe is beneficial only if its accuracy exceeds that of a text-only classifier: otherwise, it offers no advantage over a text-only classifier, which costs less to train and does not require access to hidden states.
